\section{Consequences for differentiable learning}\label{sec:resources}
A polynomial-time SQ learner can be simulated by gradient descent on an engineered differentiable model. The simulation charges its computation in the model size. This gives a consequence for the title question of \citet{FKS26}, while keeping the prescribed online-SGD problem separate.

\subsection{An engineered gradient-descent corollary}
The model used by \citet[Theorem~3d]{AKMSS21} is a directed acyclic neural network with $n+1$ inputs, including a constant input, one output, and one trainable weight per edge. Its activation is the piecewise-linear ``two-stage ramp''
\[
 \sigma(z)=\begin{cases}
 -2,&z<-3,\\ z+1,&-3\le z\le-1,\\0,&-1<z<0,\\z,&0\le z\le2,\\2,&z>2.
 \end{cases}
\]
Initialization is a specified polynomial-time computable map $W:\{0,1\}^r\to\R^p$, rather than an independent Gaussian law. Labels lie in $\{0,1\}$ and the loss is $\tfrac12(f(x)-y)^2$. Full-batch updates reuse one iid sample $S$ of size $m_s$, with a fixed constant step size and
\begin{equation}\label{eq:engineered-gradient}
 \theta_{t+1}=\theta_t-\gamma g_t,\qquad
 g_t\in\lambda\mathbb Z^p\cap[-1,1]^p,\qquad
 \left\|g_t-\frac1{m_s}\sum_{(x,y)\in S}
       [\nabla_\theta\tfrac12(f_{\theta_t}(x)-y)^2]_1\right\|_\infty
 \le\frac{3\lambda}{4}.
\end{equation}
Here $[\cdot]_1$ denotes coordinatewise clipping to $[-1,1]$. The guarantee takes the expectation over the sample and initialization after the supremum over all admissible rounded gradients. The constructed trajectories are differentiable and have gradient coordinates bounded by one, so clipping leaves them unchanged.

The imported simulation states that an SQ algorithm with $k$ queries, tolerance $\tau_0$, $r$ random bits and running time $\mathrm{TIME}$ has such a simulation with
\begin{equation}\label{eq:engineered-parameters}
 \lambda=\tau_0/16,\quad T'=2k,\quad
 m_s\lambda^2>C\bigl(k\log(1/\lambda)+\log(1/\delta)\bigr),\quad
 p=\operatorname{poly}(k,1/\tau_0,r,\mathrm{TIME},1/\delta),
\end{equation}
whose expected square loss exceeds the SQ algorithm's by at most $\delta$. The constant $C$ is universal. These are the parameters of the full-batch part of that theorem; the sample is not refreshed at each step.

\begin{corollary}[Engineered differentiable learning with high dimension]\label{cor:engineered}
Assume the PRF hypothesis of \cref{thm:prf}. Fix $c\ge1$ and $0<\epsilon<1/4$. For every sufficiently large admissible encoding length $n$ and every key $s$, the class $H_{A_s}$, encoded on $\{0,1\}^n$, has an engineered neural network and initialization of the form above, each describable and evaluable in $\operatorname{poly}_{c,\epsilon}(n)$ time, with $\operatorname{poly}_{c,\epsilon}(n)$ parameters. After $O(n^2)$ full-batch updates using $\operatorname{poly}_{c,\epsilon}(n)$ examples and the finite-precision gradient rule \eqref{eq:engineered-gradient}, thresholding the final output at $1/2$ has expected classification error at most $\epsilon$ for every marginal and every target in $H_{A_s}$. For all but a negligible fraction of keys,
\[
 \dc(H_{A_s})\ge\frac{n^c}{30\log_2 n}.
\]
The model and initialization may depend on $s$, but are fixed before the marginal and target.
\end{corollary}
\begin{proof}
Use the deterministic improper order-median learner of \cref{thm:fast} at accuracy $\epsilon/8$. It works for every key with $k=O(n^2)$ queries, tolerance $\Omega(\epsilon)$, and polynomial running time including query evaluation. Recode its labels and predictions from $\{-1,+1\}$ to $\{0,1\}$. Its square loss is half its classification loss, hence at most $\epsilon/16$ against every valid response sequence. A dyadic tolerance $\tau_0$ within a factor two of the advertised one preserves this guarantee and the asymptotic bounds.

Apply \eqref{eq:engineered-parameters} with $r=0$ and $\delta=\epsilon/16$. The sample size and model size are polynomial in $n$ for fixed $c,\epsilon$, and the expected final square loss is at most $\epsilon/8$. A thresholding mistake implies $|f(x)-y|\ge1/2$ and therefore square loss at least $1/8$. The expected classification loss is at most $\epsilon$. The PRF sign-rank conclusion is \cref{thm:prf}; it is independent of the simulation. The deterministic source learner satisfies the simulation's supremum-over-responses premise without exchanging a supremum and an expectation.
\end{proof}

The activation, architecture and initialization are engineered, and model size charges the simulated computation. In particular, the polynomial degree in the size bound can depend on $c$. The corollary establishes gradient-based distribution-independent learning of the class; it does not assert a superpolynomial separation between dimension and the size of this particular network.

\subsection{The prescribed online-SGD question}
For comparison, the remaining question fixes $X=\{-1,+1\}^n$ and a fully connected bias-free ReLU network
\[
 f_\theta(x)=\theta_L\sigma\bigl(\theta_{L-1}\sigma(\cdots\sigma(\theta_1x))\bigr),
 \qquad S=\sum_{i=1}^L n_i n_{i-1}.
\]
Each layer-$i$ weight is initialized independently with law $N(0,1/n_{i-1})$. One fixed step size is used for $T$ online updates on fresh examples with logistic loss $\log(1+e^{-yf_\theta(x)})$. The output is the sign of the sum of the final half of the iterates. The architecture, step size and horizon are fixed for the class before the marginal and target. Open Question 1 asks whether a common expected-error guarantee at $\epsilon<1/4$ forces $\dc(H)\le CTS$, or a polynomial bound in $(T,S)$. This question remains open here.

A population-gradient method with $S$ coordinates, $T$ calls, per-example coordinate bound $B_g$ and adversarial Euclidean tolerance $\delta$ has an SQ simulation with
\[
 m=TS,\qquad\tau=\delta/(B_g\sqrt S).
\]
Query each gradient coordinate divided by $B_g$. After rescaling, each error is at most $\delta/\sqrt S$, so the total squared error is at most $\delta^2$. Induction over the calls proves the simulation. Conversely, the adaptive objective $J_q(t;x,y)=tq(x,y)$ has derivative $\E q$. These observations concern gradient-oracle models. A sampled online gradient need not lie in a prescribed small ball about the population gradient, and its sampling-law guarantee does not imply success against every perturbation in that ball; see also \citet{FGV21,AKMSS21}.

\begin{proposition}[A bias-free ReLU representation]\label{thm:network}
Pad the grid domain to a cube of dimension $n=\lceil\log_2(2N^3)\rceil$, giving every padding point label $+1$. Every member is represented exactly by a bias-free ReLU network with widths $(n,2n,N+1,1)$ and at most
\[
 S=2n^2+2n(N+1)+(N+1)
\]
weights. Its common dimension complexity is at most $2N+1$. This proposition is an upper bound, not a claim that $S$ parameters are necessary.
\end{proposition}

\noindent Proof: see \cref{proof:network}.\par

A fixed $b$-bit-per-weight architecture has at most $2^{bS}$ parameter choices. Representing all $2^N$ independent labelings of a fixed line in such an architecture would require $bS\ge N$. Unrestricted real weights have no fixed bit budget, so this finite-precision count does not give a lower bound for every real-weight representation of an individual row.

\subsection{Hidden keys and class information}\label{sec:secret}
Fix the full-length template line $\ell_0:y=x+1$ and its uniform point distribution. \Cref{prop:unseen} proves a sample lower bound for a single learner independent of a random table, even when that learner is told $\ell_0$ and the marginal. The missing information is the unobserved label pattern. The computational analogue concerns a hidden PRF key.

\begin{proposition}[Hidden-key computational hardness]\label{prop:hiddenkey}
Under the PRF assumption of Theorem~\ref{thm:prf}, no probabilistic polynomial-time learner independent of $s$, with a polynomial-time evaluable output and only polynomially many labeled samples from the uniform distribution on $\ell_0$, has average error at most $1/2-\gamma(n)$ for uniform $s$, where $\gamma(n)\ge1/\operatorname{poly}(n)$. More precisely, a learner with $q(n)$ samples and advantage $\gamma(n)$ gives a PRF distinguisher of advantage at least
\[
 \gamma(n)-\frac{q(n)}{2N}.
\]
\end{proposition}

\noindent Proof: see \cref{proof:hiddenkey}.\par

A guarantee for every key, or for all but a negligible fraction of keys, implies the average guarantee ruled out by \cref{prop:hiddenkey}. The original online-SGD formulation permits architecture, step size and horizon to depend on the class. Gaussian initialization alone does not remove that dependence. Thus the hidden-key reduction applies to key-independent finite-time procedures, rather than to every class-dependent design allowed in Open Question 1. A negative answer to that question requires a small compatible network and an analysis of the specified dynamics under every marginal.
